% Chapter 8 -- Menu Commands (Part II, Reference; AN-14 menu-command
% entries: one level-1 head per menu in menu-bar order, a capture of the
% pulled-down menu, one bulleted paragraph per command with the keyboard
% equivalent).  CP-10: every menu and every item of every shipped menu of
% the desktop -- the upper bar (os/flair/shell_menus.h x kmain.c
% fd_sys_titles), the Apple menu (os/flair/apple_menu.c,
% spec/apple_menu_items.h), the Finder (os/flair/finder_menu.c, the
% predicate map; os/flair/finder_cmd.c and finder_windows.c for what each
% command does), HELLO (menu_canon.h titles, kmain.c items) and NOTES
% (kmain.c flair_ten_notes_*).  Every entry is registered (\MenuItemEntry,
% initech-flair.sty) and compared with harness/pubs/flair_inventory.c by
% test-pubs-flair in both directions; the State argument is the shipped
% enable byte (on / off / cond = lit by a Finder predicate).  The menus of
% the disk applications are in the FLAIR Accessories Guide and the
% applications' own books.
% Items that are enabled but act on nothing in this build (the upper bar's
% Quit -- shell_menus.h "a STATED, NAMED EXCEPTION ... a no-op today";
% every item of HELLO's and NOTES' bars -- band-2 choices reach only a disk
% application, kmain.c flair_live_tenant_menu_result) carry the DRAFT
% marker (D-14).
% GATES: test-finder-menu test-finder-cmd test-menu-handlers test-flair-finder-cmds test-flair-disk-windows test-flair-empty-trash test-flair-restart test-flair-shutdown test-flair-exit-to-dos test-apple-menu test-flair-apple-menu test-flair-menu test-flair-appswitch test-pubs-flair test-pubs-flair-screens
\Gates{test-finder-menu, test-finder-cmd, test-menu-handlers,
  test-flair-finder-cmds, test-flair-disk-windows, test-flair-empty-trash,
  test-flair-restart, test-flair-shutdown, test-flair-exit-to-dos,
  test-apple-menu, test-flair-apple-menu, test-flair-menu,
  test-flair-appswitch, test-pubs-flair, test-pubs-flair-screens}
\chapter{Menu Commands}
\label{ch:menus}

This chapter describes every command in the menus of the FLAIR desktop:
the upper menu bar, the Apple menu, the Finder's menus, and the menus of
HELLO and NOTES.  The menus are described in the order they appear in the
menu bars, and the commands in the order they appear in each menu.  For
each command the chapter tells you what it does, when it is available, and
its keyboard equivalent, if it has one.

A command drawn dimmed cannot be chosen.  Some commands are dimmed only
until something makes them available, such as a selected folder or a
window in front; the entry says so.  Others are dimmed at all times in this
version of FLAIR.

The menus of InitechText, the Calculator, InitechCalc, InitechPaint and the
sample applications are described in the \doctitle{FLAIR Accessories
Guide}; Initech 123 and InitechWord are described in their own books.

\section{The upper menu bar}
\index{menu bars!upper}
The upper menu bar is FLAIR's own and never changes.  At its left is the
Apple menu, described next.  Its four other menus, \menuname{File},
\menuname{Edit}, \menuname{View} and \menuname{Special}, each hold the same
two commands.

\FlairScreen[scale=4, label=fig:m-band1]{band1file}{The File menu of the
upper menu bar}

\subsection{File}
\begin{MenuReference}{SYSTEM}{File}
\MenuItemEntry{About}{}{off}{Dimmed.}
\MenuItemEntry[choosing the upper bar's Quit has no effect in this version
(the shell menus)]{Quit}{Q}{on}{Drawn available; the Finder's commands are in
the lower menu bar.}
\end{MenuReference}
\subsection{Edit}
\begin{MenuReference}{SYSTEM}{Edit}
\MenuItemEntry{About}{}{off}{Dimmed.}
\MenuItemEntry[as File > Quit]{Quit}{Q}{on}{Drawn available.}
\end{MenuReference}
\subsection{View}
\begin{MenuReference}{SYSTEM}{View}
\MenuItemEntry{About}{}{off}{Dimmed.}
\MenuItemEntry[as File > Quit]{Quit}{Q}{on}{Drawn available.}
\end{MenuReference}
\subsection{Special}
\begin{MenuReference}{SYSTEM}{Special}
\MenuItemEntry{About}{}{off}{Dimmed.}
\MenuItemEntry[as File > Quit]{Quit}{Q}{on}{Drawn available.}
\end{MenuReference}
The keyboard equivalent \CmdKey{Q} shown in these menus is not given to
the upper bar: a \key{Ctrl} key combination always goes to the menus of
the active application, in the lower bar.

\section{The Apple menu}
\index{Apple menu}
The Apple menu is at the left of the upper menu bar, and also at the left
of the lower bar while the Finder or a disk application with an Apple menu
is active; it is the same menu in every place.  FLAIR builds it again from
the data disk each time you pull it down.

\FlairScreen[scale=4, label=fig:m-apple]{applemenu}{The Apple menu}

\begin{MenuReference}{APPLE}{Apple}
\MenuItemEntry{About This Computer...}{}{on}{Shows an alert with the
version of InitechDOS, the amount of memory built into the computer, and
the largest block of that memory not in use.  Click \texttt{OK} or press
\key{Enter} to put it away.  See Figure~\ref{fig:about}.}
\MenuItemEntry{Calculator}{}{on}{Opens the Calculator accessory, the program
\file{CALCULAT.EXE} in the \file{APPS} folder.  The command is in the menu only
while that program is on the data disk; without it, the menu holds only
\menuitem{About This Computer...}.  The Calculator does not open while
another application from the disk is open.}
\end{MenuReference}

\section{The Finder's File menu}
\index{File menu (Finder)}
The Finder's menus are in the lower menu bar while the Finder is the active
application: after you click the desktop or a Finder window.

\FlairScreen[scale=4, label=fig:m-file]{filemenu}{The Finder's File menu}

Figure~\ref{fig:m-file} shows the menu with no Finder window open;
Figure~\ref{fig:kfilesel}, in Chapter~\ref{ch:disks}, shows it with a folder
selected in a Finder window, when \menuitem{Open} and \menuitem{Close
Window} are available.

\begin{MenuReference}{FINDER}{File}
\MenuItemEntry{New Folder}{N}{on}{Makes a new, empty folder in the folder
of the active Finder window, named \file{NEWFOLD} (or \file{NEWFOL02} to
\file{NEWFOL99} if that name is taken).  With no Finder window in front it
makes nothing.}
\MenuItemEntry{Open}{O}{cond}{Opens the selected folder in a window.
Available only when a folder is selected in the active Finder window.  To
open an application, double-click its icon.}
\MenuItemEntry{Print}{}{off}{Dimmed.  FLAIR does not print.}
\MenuItemEntry{Close Window}{W}{cond}{Closes the active Finder window.
Available only when a Finder window is in front.}
\MenuItemEntry{Get Info}{I}{off}{Dimmed.}
\MenuItemEntry{Duplicate}{D}{off}{Dimmed.}
\MenuItemEntry{Make Alias}{}{off}{Dimmed.}
\MenuItemEntry{Put Away}{Y}{off}{Dimmed.  To take an item out of the
Trash, drag it out of the Trash window (Chapter~\ref{ch:trash}).}
\MenuItemEntry{Find}{F}{off}{Dimmed.}
\MenuItemEntry{Page Setup}{}{off}{Dimmed.}
\MenuItemEntry{Sharing}{}{off}{Dimmed.}
\MenuItemEntry{Eject}{}{off}{Dimmed.  The data disk stays in place while
FLAIR runs.}
\end{MenuReference}

\section{The Finder's Edit menu}
\index{Edit menu (Finder)}
\FlairScreen[scale=4, label=fig:m-edit]{editmenu}{The Finder's Edit menu}

\begin{MenuReference}{FINDER}{Edit}
\MenuItemEntry{Undo}{Z}{off}{Dimmed.}
\MenuItemEntry{Cut}{X}{off}{Dimmed.}
\MenuItemEntry{Copy}{C}{off}{Dimmed.}
\MenuItemEntry{Paste}{V}{off}{Dimmed.}
\MenuItemEntry{Clear}{}{off}{Dimmed.}
\MenuItemEntry{Select All}{A}{on}{Selects every icon in the active Finder
window or, with no Finder window in front, every icon on the desktop.}
\MenuItemEntry{Show Clipboard}{}{off}{Dimmed.}
\MenuItemEntry{Preferences}{}{off}{Dimmed.}
\end{MenuReference}

\section{The Finder's View menu}
\index{View menu (Finder)}
The View menu acts on the active Finder window.  Its commands are
available only while a Finder window is in front.

\FlairScreen[scale=4, label=fig:m-view]{viewmenu}{The Finder's View menu}

\begin{MenuReference}{FINDER}{View}
\MenuItemEntry{by Icons}{}{cond}{Shows the window's contents as icons.  The
check mark beside it shows that icons are the way the window is viewed;
it is the only way in this version of FLAIR.}
\MenuItemEntry{by Small Icon}{}{off}{Dimmed.}
\MenuItemEntry{by Name}{}{off}{Dimmed.}
\MenuItemEntry{by Size}{}{off}{Dimmed.}
\MenuItemEntry{by Kind}{}{off}{Dimmed.}
\MenuItemEntry{by Label}{}{off}{Dimmed.}
\MenuItemEntry{by Date}{}{off}{Dimmed.}
\MenuItemEntry{as Buttons}{}{off}{Dimmed.}
\MenuItemEntry{as Pop-up Window}{}{off}{Dimmed.}
\MenuItemEntry{Clean Up}{}{cond}{Moves every icon of the active window to
a place of the window's grid, so that the icons stand in tidy rows.}
\MenuItemEntry{Arrange (by Name)}{}{cond}{Lays the icons of the active
window out in rows in alphabetical order of their names.}
\end{MenuReference}

\section{The Finder's Special menu}
\index{Special menu (Finder)}
\FlairScreen[scale=4, label=fig:m-special]{specialmenu}{The Finder's Special
menu, with no Finder window open and the Trash empty}

\FlairScreen[scale=4, label=fig:m-specialfull]{specialfull}{The Special
menu with a Finder window in front and something in the Trash}

\begin{MenuReference}{FINDER}{Special}
\MenuItemEntry{Clean Up}{}{cond}{The same command as \menuitem{Clean Up} in
the View menu.}
\MenuItemEntry{Empty Trash}{}{cond}{Deletes everything in the Trash from the
disk, after an alert that tells you how much it holds and asks you to
confirm (Chapter~\ref{ch:trash}).  Available only while the Trash holds something.}
\MenuItemEntry{Erase Disk}{}{off}{Dimmed.}
\MenuItemEntry{Restart}{}{on}{Asks an open application to quit, saves the
arrangement of the desktop and starts the computer again.  Restart does not
ask you to confirm.}
\MenuItemEntry{Shut Down}{}{on}{Asks an open application to quit, saves the
arrangement of the desktop and shows the message that it is safe to switch
off the computer, with a \texttt{Restart} button (Chapter~\ref{ch:starting}).}
\MenuItemEntry{Sleep}{}{off}{Dimmed.}
\MenuItemEntry{Exit to InitechDOS}{}{on}{After an alert that asks you to
confirm, asks an open application to quit, saves the arrangement of the
desktop and leaves FLAIR for the InitechDOS command prompt (Chapter~\ref{ch:starting}).
To return to FLAIR, restart the computer.}
\end{MenuReference}
If an open application declines to quit --- because you click
\texttt{Cancel} in its alert about unsaved changes --- \menuitem{Restart},
\menuitem{Shut Down} and \menuitem{Exit to InitechDOS} stop, and nothing
else happens.

\section{The Finder's Help menu}
\index{Help menu (Finder)}
\FlairScreen[scale=4, label=fig:m-help]{helpmenu}{The Finder's Help menu}

\begin{MenuReference}{FINDER}{Help}
\MenuItemEntry{About Help}{}{off}{Dimmed.}
\MenuItemEntry{Show Balloons}{}{off}{Dimmed.}
\end{MenuReference}

\section{HELLO's menus}
\index{HELLO!menus}
HELLO's menu bar is in the lower bar while the HELLO window is in front,
as it is when FLAIR starts.  It has eight menus, \menuname{File},
\menuname{Edit}, \menuname{Image}, \menuname{Layer}, \menuname{Select},
\menuname{View}, \menuname{Window} and \menuname{Help}, and no Apple menu.
Each menu holds the same two commands.

\FlairScreen[scale=4, label=fig:m-hello]{hellofile}{The File menu of
HELLO's menu bar}

\newcommand\HelloMenu[1]{%
\subsection{#1}
\begin{MenuReference}{HELLO}{#1}
\MenuItemEntry[HELLO's menu commands act on nothing in this version: a
choice in the lower bar reaches only an application from the disk
(kmain.c)]{New}{N}{on}{Drawn available.}
\MenuItemEntry[as New]{Open}{O}{on}{Drawn available.}
\end{MenuReference}}
\HelloMenu{File}
\HelloMenu{Edit}
\HelloMenu{Image}
\HelloMenu{Layer}
\HelloMenu{Select}
\HelloMenu{View}
\HelloMenu{Window}
\HelloMenu{Help}

\section{NOTES' menus}
\index{NOTES!menus}
NOTES' menu bar is in the lower bar while the NOTES window is in front.
It has three menus, \menuname{File}, \menuname{Edit} and
\menuname{Notes}, and no Apple menu.

\FlairScreen[scale=4, label=fig:m-notesfile]{notesfile}{NOTES' File menu}

\subsection{File}
\begin{MenuReference}{NOTES}{File}
\MenuItemEntry[NOTES' menu commands act on nothing in this version: a
choice in the lower bar reaches only an application from the disk
(kmain.c)]{New}{N}{on}{Drawn available.}
\MenuItemEntry[as New]{Open}{O}{on}{Drawn available.}
\MenuItemEntry[as New]{Save}{S}{on}{Drawn available.}
\MenuItemEntry[as New]{Quit}{Q}{on}{Drawn available.  To close NOTES, click
its close box.}
\end{MenuReference}

\FlairScreen[scale=4, label=fig:m-notesedit]{notesedit}{NOTES' Edit menu}

\subsection{Edit}
\begin{MenuReference}{NOTES}{Edit}
\MenuItemEntry[as File > New]{Undo}{Z}{on}{Drawn available.}
\MenuItemEntry[as File > New]{Cut}{X}{on}{Drawn available.}
\MenuItemEntry[as File > New]{Copy}{C}{on}{Drawn available.}
\MenuItemEntry[as File > New]{Paste}{V}{on}{Drawn available.}
\end{MenuReference}

\FlairScreen[scale=4, label=fig:m-notesnotes]{notesnotes}{NOTES' Notes menu}

\subsection{Notes}
\begin{MenuReference}{NOTES}{Notes}
\MenuItemEntry[as File > New]{About Notes}{}{on}{Drawn available.}
\end{MenuReference}
The keyboard equivalents shown in the menus of HELLO and NOTES are not
given to them: while HELLO or NOTES is in front, a \key{Ctrl} key
combination does nothing.
